\section{Introduction}

\subsection{Background and motivation}

Every US public company must publish its financial results. Each quarter it files a report called a
10-Q with the Securities and Exchange Commission (SEC), and each year a longer annual report, the
10-K. These filings answer questions such as ``What were 3M's net sales in the second quarter of
2023?'' or ``Why did Honeywell's sales increase?''. They are free to download from the SEC's EDGAR
website~\cite{sec-edgar}, but they are long: a single 10-Q often has 50 pages, and more than 140,000
10-Qs have been filed since 2019.

Large language models (LLMs) such as ChatGPT or Claude can answer such questions in fluent text, but
they make two kinds of mistakes that matter in finance. They may state a number from memory that is
wrong or out of date, and they often get arithmetic wrong even when the correct figures are in front
of them~\cite{financebench,finqa}. A wrong number given with confidence is worse than no answer.

This project builds a question-answering system for company filings that \emph{cites the source of
every number and every reason}, so that a reader can check each claim. The central idea is to give
the LLM as little responsibility as possible: it only interprets the question and phrases the final
answer, while ordinary code finds the figures, computes changes and checks the citations. To answer
questions about any company, the system needs the data of every filer, so the project also has a
cloud and big data side: tens of gigabytes of SEC data are processed with Spark on AWS.

\subsection{Literature review}

\paragraph{Financial question answering.} Questions about financial reports mix text, tables and
arithmetic. The FinQA~\cite{finqa} and TAT-QA~\cite{tatqa} benchmarks pair questions with a table and
some text from an annual report; both found that models are much better at finding the right numbers
than at computing with them. FinanceBench~\cite{financebench} is closer to real use: 150 questions
written by financial analysts about real 10-K and 10-Q filings. Its authors found that a standard
LLM pipeline answered most of them incorrectly or not at all. We use it as an external test set.

\paragraph{Retrieval-augmented generation (RAG).} RAG~\cite{rag} is the standard way to let an LLM
answer from documents it was not trained on: a search step selects the most relevant passages and
gives them to the LLM together with the question. Search can be \emph{lexical}, such as
BM25~\cite{bm25}, which scores passages by the words they share with the question, or
\emph{semantic}, which turns texts into vectors (embeddings) with a model such as
MiniLM~\cite{minilm,sbert} and compares how similar they are. Reciprocal rank fusion~\cite{rrf}
combines the two rankings. RAG does not stop the LLM from adding facts that are not in the
passages~\cite{hallucination}, and it does not help with arithmetic. Program-of-thought
prompting~\cite{pot} moves arithmetic out of the LLM into code; we go one step further and take the
most common numbers directly from structured data.

\paragraph{Big data and small models.} Apache Spark~\cite{spark} processes large datasets in
parallel, on one machine or on a cluster. Postgres with the pgvector extension~\cite{pgvector} stores
embeddings next to ordinary tables, so figures, filings and passages can live in one database. Small
open models such as Qwen2.5~\cite{qwen25} run on a laptop at no cost per question but are less
reliable. Whether a careful design makes such models useful is one of our research questions.

\subsection{Objectives}

\begin{enumerate}
  \item Build a question-answering workflow over SEC filings in which code, not the LLM, finds and
        computes every number, and every claim cites a filing, a statement line or a passage.
  \item Build a data pipeline on AWS that uses Spark to load the figures and the 10-Q text of every
        SEC filer since 2019 into one database.
  \item Compare the workflow with simpler systems on our own question set and on FinanceBench, using
        a hosted LLM (Claude Haiku 5.5) and small local LLMs (Qwen2.5 with 3 and 7 billion
        parameters).
  \item Measure the cloud side: time per step, throughput, Spark scaling and the time needed to add
        new filings.
\end{enumerate}
